%% Chapter 5 — Readout Architecture
\section{Readout Architecture}
\label{sec:readout}

\subsection{DDS-Based Signal Generation}

\QICK{} avoids the analogue mixing of a standard heterodyne readout chain, in
which a local oscillator (LO) up-converts a baseband envelope to the resonator
frequency $f_r$ at the price of image products, LO leakage, and IQ imbalance
correction.  Both the readout drive and the
demodulation reference are instead produced by the on-chip DDS engines, which operate at
the direct output frequency of the RFSoC converters (up to several GHz).  The
waveform generated for qubit readout has the form
\begin{equation}
  s(t) = A \cdot \text{env}(t) \cdot \cos(2\pi f_r t + \phi_\text{res}),
\end{equation}
where $\text{env}(t)$ is the pulse envelope (constant, Gaussian, Derivative
Removal by Adiabatic Gate (DRAG), etc.),
$f_r$ is programmed directly into the DDS frequency register, and
$\phi_\text{res}$ is the phase correction from calibration.  The residual
imperfections are digital: quantisation and DDS phase-truncation spurs, which are
deterministic and can be mitigated by appropriate filtering and sample-rate
choice, plus random clock jitter.

\subsection{Decimated vs. Accumulated Readout}
\label{subsec:decimated_vs_accum}

After the ADC samples the incoming signal, the \QICK{} readout block demodulates
it by multiplying with the NCO reference tone and low-pass filtering.  Two output
modes are available:

\paragraph{Decimated readout.}
Stores the full time-domain I/Q waveform over the readout window (e.g.\ 1000
samples).  This is essential for setup: it reveals when the readout pulse arrives
relative to the trigger (\texttt{adc\_trig\_offset}), whether the integration
window is correctly positioned, and whether the signal shape is as expected.  The
per-sample noise on the decimated output is $\sim$\SI{1.9}{ADU} RMS; this is the
raw sample-level noise floor, distinct from the per-shot integrated IQ noise
reported for accumulated readout below.

\paragraph{Accumulated readout.}
The FPGA sums all ADC samples within the readout window into a single $(I, Q)$
pair in hardware, one pair per shot.  This is the standard mode for SSRO: each
shot yields one IQ point, and the discrimination classifier operates on that
point.  The per-shot integrated IQ noise is determined by the within-window
averaging over $T_\text{win}$ samples, not by the number of repeated shots.
From the real SSRO data, the per-shot cluster widths are approximately
\SI{0.65}{}--\SI{0.97}{ADU} per quadrature (state-dependent), with an IQ
separation of \SI{2.53}{ADU} and a projected SNR of 2.46, numerically the
Mahalanobis distance between the clusters
(Section~\ref{subsec:integrated_iq}).  The single-shot
assignment fidelity extracted from these clusters is 0.906 (the LDA baseline of
Section~\ref{subsec:integrated_iq}), a typical value for dispersive readout
without a quantum-limited parametric amplifier in the chain.

For applications that do \emph{not} require single-shot resolution, such as
resonator spectroscopy or Rabi calibration, the same shot can be repeated
$N_\text{reps}$ times and the results averaged, reducing the noise on the
\emph{mean} phasor as $1/\sqrt{N_\text{reps}}$:
\begin{equation}
  \sigma_\text{mean} = \frac{\sigma_\text{single}}{\sqrt{N_\text{reps}}}.
\end{equation}
It cannot
be applied to SSRO, since averaging destroys the per-shot state information.

The practical workflow is to find the correct \texttt{adc\_trig\_offset} and
verify signal quality with \texttt{acquire\_decimated()}, then take all data with
\texttt{acquire()} (accumulated): single shots for SSRO, repeat-averaged shots for
spectroscopy and calibration.

\subsection{Resonator Spectroscopy}
\label{subsec:spectroscopy}

The resonator bare frequency $\omega_r$ and linewidth $\kappa$ are found by
sweeping the probe frequency around the estimated resonator frequency and measuring
the transmitted or reflected IQ amplitude.  In \QICK{}, there are two sweep
strategies:

\paragraph{FPGA-internal sweep (\texttt{add\_loop}).}
On the tProcessor v2 firmware a parameter such as the drive frequency can be swept
inside the FPGA with \texttt{add\_loop}, stepping through all points with no Python
round-trips.  However, the readout down-conversion frequency must track the drive to
keep the demodulated signal at baseband, and \texttt{add\_loop} cannot step it: it is
fixed per program.  For spectroscopy, where the two must move together, this sweep
alone is insufficient.

\paragraph{Python outer loop.}
A Python loop instead compiles and runs one \texttt{AveragerProgramV2} per frequency
point, setting the generator and matching readout down-conversion frequencies
together each time (\texttt{add\_readoutconfig} with \texttt{gen\_ch} lets \QICK{}
compute the ADC alias automatically).  The per-point overhead is a few milliseconds,
but it can sweep any parameter, including the readout frequency, and is the approach
used here.  At each point the magnitude $|I + jQ|$ of the accumulated phasor is
recorded; the resonator shows up as the dip (reflection) or peak (transmission) in
the swept response.  The workflow was implemented and validated in loopback, where
the swept response traces the front-end passband; with no device connected, no
resonator dip was measured.

\subsection{Multiplexed Readout}
\label{subsec:multiplex}

For multi-qubit systems, simultaneous readout of several qubits on the same
physical line requires \emph{multiplexed readout}: a broadband probe signal
containing multiple tones at the individual resonator frequencies $\{f_{r,k}\}$~\cite{Heinsoo2018,Walter2017}.
The \QICK{} framework supports this through a \emph{Polyphase Filter Bank} (PFB)
readout mode.

A PFB is a bank of digital bandpass filters implemented as a single, efficient DSP
operation (polyphase decomposition of a prototype filter followed by an FFT).  It
takes the broadband ADC stream and simultaneously demodulates $N$ equally-spaced
frequency channels, analogous to a digital spectrum analyser.  Advantages over
the single-tone DDS approach include:
\begin{itemize}
  \item No need to retune the ADC NCO between qubits.
  \item The per-sample cost of demodulating all $N$ channels scales as $O(\log N)$
        through the shared FFT, rather than the $O(N)$ of $N$ independent
        demodulators.
  \item Sidesteps the ADC-frequency-sweep limitation: each PFB output channel is
        a fixed digital filter, not a tunable DDS.
\end{itemize}

Multiplexed readout via the \QICK{} PFB firmware was not implemented here.  The
single-tone DDS architecture described in
Sections~\ref{subsec:decimated_vs_accum} and~\ref{subsec:spectroscopy} was used
throughout.  A natural next step is to extend the architecture to cover the four
resonator frequencies of one of Qilimanjaro's chips: 6.505, 6.730, 6.978,
and \SI{7.389}{\giga\hertz}.  These frequencies lie above the \SI{6}{\giga\hertz}
upper edge of the XM655 balun passband, so simultaneous multiplexed readout at all
four tones would require direct HD-header connections with discrete baluns and
operation of the RF-ADC above its specified input bandwidth, as
discussed in Section~\ref{subsec:xm655}.  The PFB channel spacing and prototype
filter bandwidth would also need to be verified against the \SI{225}{\mega\hertz}
minimum separation between adjacent resonators in this device.
